\chapter{Analiza wymagań}
\label{chap:algs}

\section{Koncepcja gry}
\section{Przebieg gry}
\section{Tryby rozgrywki}
\section{Wymagania funkcjonalne}

\section{Wymagania niefunkcjonalne (Filip Jezierski)}
Wymagania niefunkcjonalne określają kryteria jakościowe, które aplikacja powinna spełniać, aby zapewnić odpowiednie doświadczenie użytkownika oraz efektywne działanie systemu. W kontekście opracowywanej aplikacji mobilnej, kluczowe wymagania niefunkcjonalne obejmują aspekty takie jak wydajność, bezpieczeństwo oraz użyteczność.

\subsection{Wymagania jakościowe}
\begin{itemize}
    \item obsługa różnych rozdzielczości ekranów i urządzeń mobilnych,
    \item dostępność aplikacji na popularnych platformach mobilnych (Android, iOS), a także jako aplikacja webowa,
    \item intuicyjny i responsywny interfejs użytkownika,
    \item aplikacja powinna w odpowiedni sposób obsługiwać błędy i wyjątki, takie jak nieprawidłowe dane wejściowe czy brak internetu, zapewniając stabilność działania,
    \item ochrona danych użytkowników poprzez implementację odpowiednich mechanizmów bezpieczeństwa - szyfrowanie danych, bezpieczne logowanie, ochrona przed atakami typu SQL Injection i XSS.
\end{itemize}

\subsection{Wymagania wydajnościowe}
\begin{itemize}
    \item szybkie (<2 s) ładowanie aplikacji i minimalny czas oczekiwania na reakcję interfejsu użytkownika,
    \item ładowanie zdjęć lokalizacji w czasie nie dłuższym niż 3 sekundy,
    \item szybkie (<10 s) ładowanie danych geolokalizacyjnych zapewniające płynne działanie trybu gry opartego na lokalizacji,
\end{itemize}

\section{Wybór platformy (Filip Jezierski)}
Wybór odpowiedniej platformy dla opracowywanej aplikacji mobilnej jest kluczowym elementem procesu projektowego, mającym istotny wpływ na jej funkcjonalność, dostępność oraz doświadczenie użytkownika, a także łatwość wdrożenia i utrzymania. Na etapie analizy wymagań podjęto decyzję o stworzeniu dwóch testowych wersji aplikacji: aplikacji webowej typu PWA (Progressive Web App) oraz natywnej aplikacji mobilnej wykorzystującej technologię .NET MAUI (Multi-platform App UI). Obie platformy oferują unikalne korzyści i wyzwania, które zostały dokładnie rozważone w kontekście specyfiki projektu.

\subsection{Natywna aplikacja mobilna}
Aplikacje natywne są tworzone z wykorzystaniem dedykowanych narzędzi i języków programowania przeznaczonych dla konkretnych systemów operacyjnych, takich jak Android czy iOS. W przypadku technologii .NET MAUI możliwe jest tworzenie aplikacji wieloplatformowych przy zachowaniu wspólnej bazy kodu, co znacząco skraca czas implementacji i ułatwia utrzymanie.

Jednakże proces publikacji aplikacji natywnych w sklepach takich jak Google Play czy Apple App Store wiąże się z dodatkowymi wymaganiami i procedurami zatwierdzania, a także z koniecznością budowania i testowania aplikacji na różne systemy operacyjne, co może wydłużyć czas wdrożenia.

\subsection{Aplikacja webowa PWA}
Aplikacja webowa typu PWA (Progressive Web App) łączy cechy klasycznej strony internetowej z funkcjonalnością aplikacji mobilnej. Dzięki wykorzystaniu nowoczesnych technologii webowych, możliwe jest stworzenie aplikacji działającej w trybie offline, instalowanej bezpośrednio z przeglądarki oraz zapewniającej wrażenia zbliżone do aplikacji natywnej.

Ze względu na łatwość dostępu (poprzez przeglądarkę internetową) oraz brak konieczności przechodzenia przez proces zatwierdzania w sklepach aplikacji, PWA stanowi atrakcyjną opcję dla naszego projektu. Jednakże, pewne ograniczenia związane z dostępem do funkcji urządzenia (takich jak geolokalizacja) oraz wydajnością w porównaniu do aplikacji natywnych muszą być uwzględnione podczas projektowania i implementacji.

\section{Ograniczenia projektu (Filip Jezierski)}
Ograniczenia projektu wynikają przede wszystkim z przyjętych założeń technologicznych oraz charakteru opracowywanego rozwiązania. Kluczowe znaczenie mają tu specyfika aplikacji webowej typu PWA, zastosowanie mechanizmów geolokalizacji jako podstawowego elementu funkcjonalnego, a także decyzje dotyczące sposobu pozyskiwania i wykorzystania materiałów wizualnych. Wymienione czynniki w naturalny sposób wpływają na zakres możliwości implementacyjnych, wydajność rozwiązania oraz doświadczenie użytkownika końcowego.

\subsection{Środowisko aplikacji}
Aplikacja będzie udostępniana jako aplikacja webowa, dostępna przez przeglądarkę internetową, z możliwością instalacji na ekranie głównym urządzenia. Funkcjonalność trybu gry wykorzystującego geolokalizację będzie dostępna wyłącznie na urządzeniach mobilnych, które wspierają tę technologię. W przypadku urządzeń stacjonarnych, nawet jeśli geolokalizacja jest technicznie możliwa, jej dokładność jest niewystarczająca do zapewnienia poprawnego działania tej funkcji.

\subsection{Zdjęcia lokalizacji}
W projekcie zrezygnowano z wykorzystywania zdjęć 360° pochodzących z zewnętrznych źródeł, takich jak interfejsy API udostępniane przez firmy trzecie. Decyzję tę podjęto z uwagi na chęć uniknięcia potencjalnych problemów związanych z licencjonowaniem, ograniczoną dostępnością materiałów oraz zróżnicowaną jakością zdjęć. Dodatkowo korzystanie z takich źródeł mogłoby wprowadzać ograniczenia dotyczące liczby i różnorodności dostępnych lokalizacji.

W celu uproszczenia procesu pozyskiwania materiałów wizualnych oraz umożliwienia użytkownikom współtworzenia bazy lokalizacji, aplikacja będzie wykorzystywać standardowe zdjęcia dwuwymiarowe. Tego typu fotografie nie wymagają specjalistycznego sprzętu i mogą być wykonywane przy użyciu dowolnego smartfona. Użytkownicy będą mieli możliwość przesyłania własnych zdjęć lokalizacji, które po weryfikacji przez moderatorów zostaną dodane do bazy danych aplikacji.

Takie rozwiązanie zapewnia pełną kontrolę nad jakością i zgodnością materiałów z założeniami projektu, a jednocześnie pozwala na stopniowe rozszerzanie bazy lokalizacji o nowe zdjęcia, dostarczane zarówno przez zespół projektowy, jak i społeczność użytkowników.

\section{Kryteria akceptacji (Filip Jezierski)}
\textbf{Ogólne kryteria:}
\begin{itemize}
    \item Aplikacja działa na urządzeniach mobilnych z systemami Android i iOS, na komputerach stacjonarnych oraz w przeglądarkach internetowych
    \item Intuicyjny interfejs użytkownika umożliwiający łatwe korzystanie z aplikacji
    \item Responsywny design dostosowujący się do różnych rozmiarów ekranów, w szczególności urządzeń mobilnych
    \item Stabilne działanie aplikacji bez krytycznych błędów i awarii
    \item Testy manualne wszystkich zaplanowanych funkcjonalności przeprowadzone na wszystkich wspieranych platformach
\end{itemize}

\textbf{Zaimplementowane funkcjonalności:}
\begin{itemize}
    \item Rejestracja i logowanie użytkowników
    \item Wybór trybu rozgrywki i poziomu trudności
    \item Dwa tryby gry:
    \begin{itemize}
        \item odgadywanie lokalizacji na podstawie zdjęcia poprzez wskazanie jej na mapie
        \item odnajdywanie lokalizacji na podstawie zdjęcia i potwierdzanie obecności poprzez geolokalizację
    \end{itemize}
    \item Możliwość przerywania i kontynuowania rozgrywki
    \item Wyświetlanie podsumowania wyników każdej rundy gry
    \item Historia rozgrywek dostępna w profilu użytkownika
    \item Ogólnodostępny ranking graczy
    \item Dodawanie lokalizacji przez użytkowników (po weryfikacji przez moderatorów)
\end{itemize}
